% GENERATED FILE. Edit canonical lessons, then run npm run generate:books.
% canonical-source-hash: fnv1a64:8d075ee0c06454a5
% canonical-lessons: UR-C267-bad, UR-W267-ain, UR-C267-kha-kar, UR-C267-pehle-phir, UR-C267-sara-ka-safar, UR-R267-a-story-in-order

\chapter{A Story in Order}
\label{ch:ur-a-story-in-order}
\hlchaptermodality{267}

\begin{chapteropening}
\textbf{By the end of this chapter:} \emph{I can tell what happened in order (pehle, phir, us ke bād, ākhir meṅ), join two actions with kar (khā kar soyā), and read a short story.}

Complete UR-R267-a-story-in-order: tell what happened in order, joining the steps.
\end{chapteropening}

\section[\texorpdfstring{\ur{بعد} (bād)}{بعد (bād)}]{\ur{بعد} (bād) — after}

\label{lesson:UR-C267-bad}



\begin{quote}
Before the new one: say \textquotedblleft{}during the journey\textquotedblright{}. (\emph{safar ke daurān}.) Then \textquotedblleft{}before\textquotedblright{}. (\emph{pehle}.)
\end{quote}

\subsection*{You'll want to know: \ur{بعد}}
\textbf{\ur{بعد}} — \emph{bād} — \textquotedblleft{}after\textquotedblright{}.

\begin{quote}
\textbf{\ur{کھانے} \ur{کے} \ur{بعد} \ur{میں} \ur{سویا۔}}
\end{quote}

\begin{quote}
\emph{khāne ke bād maiṅ soyā.} — \textquotedblleft{}After the meal I slept.\textquotedblright{}
\end{quote}

Like \textbf{\ur{دوران}} (\emph{daurān}), it stands after \textbf{\ur{کے}} (\emph{ke}). With \textbf{\ur{میں}} (\emph{meṅ}) and no noun, it is \textquotedblleft{}later\textquotedblright{}:

\begin{quote}
\textbf{\ur{میں} \ur{بعد} \ur{میں} \ur{آؤں} \ur{گا۔}}
\end{quote}

\begin{quote}
\emph{maiṅ bād meṅ āūṅ gā.} — \textquotedblleft{}I'll come later.\textquotedblright{}
\end{quote}

\begin{cousinweb}
\textbf{\ur{بعد}} (\emph{bād}) is usually traced, through Persian, to Arabic \emph{baʿd}, \textquotedblleft{}after\textquotedblright{}. The letter in its middle is new to your hand; the next lesson takes it out.
\end{cousinweb}

\subsection*{Your turn}
Say these aloud:

\begin{itemize}
\raggedright
  \item \emph{bād}, then \emph{khāne ke bād}
  \item \textquotedblleft{}After the meal she went to the market\textquotedblright{} — \emph{khāne ke bād voh bāzār gaī}
  \item \textquotedblleft{}I'll go later\textquotedblright{} — \emph{maiṅ bād meṅ jāūṅ gā}
\end{itemize}

\begin{tcolorbox}[breakable,colback=teal!4,colframe=teal!35!black,title={Before you move on}]
What does \textbf{\ur{بعد}} mean? (\textquotedblleft{}After\textquotedblright{}.) \textquotedblleft{}After the meal\textquotedblright{}? (\textbf{\ur{کھانے} \ur{کے} \ur{بعد}}.) \textquotedblleft{}Later\textquotedblright{}? (\textbf{\ur{بعد} \ur{میں}}.)
\end{tcolorbox}
\section[\texorpdfstring{\ur{ع} (ain)}{ع (ain)}]{\ur{ع} (ain) — the letter inside \ur{بعد}}

\label{lesson:UR-W267-ain}



\begin{quote}
Before the new one: write \ur{ؤ} once, and name it. (\emph{vāʾo hamza}.)

Now say \textbf{\ur{بعد}} (\emph{bād}), \textquotedblleft{}after\textquotedblright{}. This lesson takes one letter out of it and puts it in your hand.
\end{quote}

\subsection*{Script you'll notice: \ur{ع}}
\textbf{\ur{ع}} — \emph{ain}.

It sits in the middle of \textbf{\ur{بعد}} (\emph{bād}). In Urdu it has no sound of its own: it colours the vowel beside it, so the \emph{a} of \textbf{\ur{بعد}} is said long, \emph{bād}. In the middle of a word it is a small closed loop on the line, joined on both sides. Standing alone it opens out: a small open head, then a deep bowl that sweeps below the line.

\subsection*{Writing: \ur{ع}}
Put your pen on \ur{ع} and follow its line. Copy the shape you can see — slowly, and larger than it is printed.

\begin{quote}
This book does not yet tell you \textbf{where to start this letter or which way to travel}. It is not written down here until it can be written down with a source. Copying what is in front of you needs no such source, and it is how the shape gets into your hand in the meantime.
\end{quote}

\subsection*{Your turn}
\begin{itemize}
\raggedright
  \item \emph{Look:} at \textbf{\ur{بعد}} and point at \ur{ع} inside it
  \item \emph{Trace it:} \ur{ع} three times, saying \emph{bād} as you finish each one
  \item \emph{Write it:} \textbf{\ur{بعد}} once, joined, with the small loop in its middle
  \item \emph{Read it:} \textbf{\ur{بعد}} aloud once more
\end{itemize}

\begin{tcolorbox}[breakable,colback=teal!4,colframe=teal!35!black,title={Before you move on}]
Which letter is this — \ur{ع}? (\textbf{ain}: it colours the vowel.) Which word did it come from? (\textbf{\ur{بعد}}, \emph{bād}, \textquotedblleft{}after\textquotedblright{}.)
\end{tcolorbox}
\section[\texorpdfstring{\ur{کھا} \ur{کر} (khā kar)}{کھا کر (khā kar)}]{\ur{کھا} \ur{کر} (khā kar) — one thing, then the next}

\label{lesson:UR-C267-kha-kar}



\begin{quote}
Say \textquotedblleft{}I ate\textquotedblright{} and \textquotedblleft{}I slept\textquotedblright{}. (\emph{maiṅ ne khāyā}, \emph{maiṅ soyā}.) Two events. This lesson joins them into one sentence.
\end{quote}

\begin{grammarlens}[title={Grammar Lens: the stem and \ur{کر}}]
\begin{quote}
\textbf{\ur{کھا} \ur{کر} \ur{سویا۔}}
\end{quote}

\begin{quote}
\emph{khā kar soyā.} — \textquotedblleft{}Having eaten, he slept.\textquotedblright{}
\end{quote}

The first action keeps only its stem, \textbf{\ur{کھا}} (\emph{khā}), and takes \textbf{\ur{کر}} (\emph{kar}), written apart. The last verb carries the person and the time:

\begin{quote}
\textbf{\ur{میں} \ur{کھانا} \ur{کھا} \ur{کر} \ur{سویا۔}}
\end{quote}

\begin{quote}
\emph{maiṅ khānā khā kar soyā.} — \textquotedblleft{}I ate, then slept.\textquotedblright{}
\end{quote}

\begin{quote}
\textbf{\ur{سارا} \ur{گھر} \ur{آ} \ur{کر} \ur{بیٹھی۔}}
\end{quote}

\begin{quote}
\emph{Sārā ghar ā kar baiṭhī.} — \textquotedblleft{}Sara came home and sat down.\textquotedblright{}
\end{quote}

\begin{quote}
\textbf{\ur{ہم} \ur{کھا} \ur{کر} \ur{جائیں} \ur{گے۔}}
\end{quote}

\begin{quote}
\emph{ham khā kar jāeṅ ge.} — \textquotedblleft{}We'll eat, then go.\textquotedblright{}
\end{quote}

You have said one for a long time: \textbf{\ur{مل} \ur{کر}} (\emph{mil kar}), \textquotedblleft{}having met\textquotedblright{}, in the greeting that ends with \textbf{\ur{خوشی} \ur{ہوئی}}.
\end{grammarlens}

\subsection*{You'll want to know: \ur{کر} \ur{کے}, and the last verb}
\textbf{\ur{کرنا}} (\emph{karnā}) makes \textbf{\ur{کر} \ur{کے}} (\emph{kar ke}), not a doubled \textbf{\ur{کر}}:

\begin{quote}
\textbf{\ur{کام} \ur{کر} \ur{کے} \ur{میں} \ur{گھر} \ur{آیا۔}}
\end{quote}

\begin{quote}
\emph{kām kar ke maiṅ ghar āyā.} — \textquotedblleft{}I finished the work and came home.\textquotedblright{}
\end{quote}

The last verb decides \textbf{\ur{نے}} (\emph{ne}). With \textbf{\ur{پی}} (\emph{pī}), \textquotedblleft{}drank\textquotedblright{}, the doer takes it:

\begin{quote}
\textbf{\ur{اس} \ur{نے} \ur{کھا} \ur{کر} \ur{چائے} \ur{پی۔}}
\end{quote}

\begin{quote}
\emph{us ne khā kar chāy pī.} — \textquotedblleft{}He ate, then drank tea.\textquotedblright{}
\end{quote}

\subsection*{Your turn}
Say these aloud:

\begin{itemize}
\raggedright
  \item \textquotedblleft{}I ate, then slept\textquotedblright{} — \emph{maiṅ khā kar soyā}
  \item \textquotedblleft{}She ate, then slept\textquotedblright{} — \emph{voh khā kar soī}
  \item \textquotedblleft{}We ate, then slept\textquotedblright{} — \emph{ham khā kar soe}
\end{itemize}

\begin{tcolorbox}[breakable,colback=teal!4,colframe=teal!35!black,title={Before you move on}]
\textquotedblleft{}Having eaten\textquotedblright{}? (\textbf{\ur{کھا} \ur{کر}}.) \textquotedblleft{}Having done the work\textquotedblright{}? (\textbf{\ur{کام} \ur{کر} \ur{کے}}.) Which verb decides \textbf{\ur{نے}}? (The last one.)
\end{tcolorbox}
\section[\texorpdfstring{\ur{پہلے،} \ur{پھر} (pehle, phir)}{پہلے، پھر (pehle, phir)}]{\ur{پہلے،} \ur{پھر} (pehle, phir) — first, then, after that, in the end}

\label{lesson:UR-C267-pehle-phir}



\begin{quote}
Say \textquotedblleft{}having eaten\textquotedblright{}. (\emph{khā kar}.) Then \textquotedblleft{}before\textquotedblright{} and \textquotedblleft{}then\textquotedblright{}. (\emph{pehle}, \emph{phir}.) A story puts its steps in a row; these words mark the row.
\end{quote}

\begin{grammarlens}[title={Grammar Lens: first and then}]
\begin{quote}
\textbf{\ur{پہلے} \ur{میں} \ur{نہایا۔}}
\end{quote}

\begin{quote}
\emph{pehle maiṅ nahāyā.} — \textquotedblleft{}First I had a bath.\textquotedblright{}
\end{quote}

\begin{quote}
\textbf{\ur{پھر} \ur{میں} \ur{نے} \ur{کھانا} \ur{کھایا۔}}
\end{quote}

\begin{quote}
\emph{phir maiṅ ne khānā khāyā.} — \textquotedblleft{}Then I ate.\textquotedblright{}
\end{quote}

\textbf{\ur{پہلے}} (\emph{pehle}) opens the first step and \textbf{\ur{پھر}} (\emph{phir}) the next. Each goes at the front of its line, so the order is clear before the verb arrives.
\end{grammarlens}

\subsection*{You'll want to know: after that, in the end}
\begin{quote}
\textbf{\ur{اس} \ur{کے} \ur{بعد} \ur{میں} \ur{دفتر} \ur{گیا۔}}
\end{quote}

\begin{quote}
\emph{us ke bād maiṅ daftar gayā.} — \textquotedblleft{}After that I went to the office.\textquotedblright{}
\end{quote}

\begin{quote}
\textbf{\ur{آخر} \ur{میں} \ur{سارا} \ur{گھر} \ur{آئی۔}}
\end{quote}

\begin{quote}
\emph{ākhir meṅ Sārā ghar āī.} — \textquotedblleft{}In the end Sara came home.\textquotedblright{}
\end{quote}

\textbf{\ur{اس} \ur{کے} \ur{بعد}} (\emph{us ke bād}) is \textquotedblleft{}after that\textquotedblright{}: \textbf{\ur{بعد}} with \textbf{\ur{اس}} (\emph{us}). \textbf{\ur{آخر} \ur{میں}} (\emph{ākhir meṅ}) closes the story.

\subsection*{Your turn}
Say these aloud:

\begin{itemize}
\raggedright
  \item \textquotedblleft{}First I went to the market\textquotedblright{} — \emph{pehle maiṅ bāzār gayā}
  \item \textquotedblleft{}Then I went to the market\textquotedblright{} — \emph{phir maiṅ bāzār gayā}
  \item \textquotedblleft{}In the end I went to the market\textquotedblright{} — \emph{ākhir meṅ maiṅ bāzār gayā}
\end{itemize}

\begin{tcolorbox}[breakable,colback=teal!4,colframe=teal!35!black,title={Before you move on}]
\textquotedblleft{}First\textquotedblright{}? (\textbf{\ur{پہلے}}.) \textquotedblleft{}After that\textquotedblright{}? (\textbf{\ur{اس} \ur{کے} \ur{بعد}}.) \textquotedblleft{}In the end\textquotedblright{}? (\textbf{\ur{آخر} \ur{میں}}.)
\end{tcolorbox}
\section[(reading)]{(reading) — Sara's trip to Lahore}

\label{lesson:UR-C267-sara-ka-safar}



\begin{quote}
Ninety words of Urdu in a row, every one met before: Sara's trip, in order.
\end{quote}

\begin{tcolorbox}[breakable,skin=enhanced,colback=orange!6,colframe=orange!45!black,boxrule=0.5pt,arc=1mm,left=6pt,right=6pt,top=4pt,bottom=4pt,fonttitle=\bfseries,title={Read this}]
Line by line, with the romanization under each:

\begin{quote}
\textbf{\ur{پچھلے} \ur{ہفتے} \ur{سارا} \ur{لاہور} \ur{گئی۔}}
\end{quote}

\begin{quote}
\emph{pichhle hafte Sārā Lāhaur gaī.}
\end{quote}

\begin{quote}
\textbf{\ur{اس} \ur{کی} \ur{دوست} \ur{لاہور} \ur{میں} \ur{رہتی} \ur{ہے۔}}
\end{quote}

\begin{quote}
\emph{us kī dost Lāhaur meṅ rahtī hai.}
\end{quote}

\begin{quote}
\textbf{\ur{پہلے} \ur{اس} \ur{نے} \ur{اسٹیشن} \ur{پر} \ur{ٹکٹ} \ur{لیا۔}}
\end{quote}

\begin{quote}
\emph{pehle us ne isṭeshan par ṭikaṭ liyā.}
\end{quote}

\begin{quote}
\textbf{\ur{پھر} \ur{وہ} \ur{ریل} \ur{گاڑی} \ur{میں} \ur{بیٹھی۔}}
\end{quote}

\begin{quote}
\emph{phir voh rel gāṛī meṅ baiṭhī.}
\end{quote}

\begin{quote}
\textbf{\ur{سفر} \ur{کے} \ur{دوران} \ur{وہ} \ur{کتاب} \ur{پڑھ} \ur{رہی} \ur{تھی۔}}
\end{quote}

\begin{quote}
\emph{safar ke daurān voh kitāb paṛh rahī thī.}
\end{quote}

\begin{quote}
\textbf{\ur{باہر} \ur{بارش} \ur{ہو} \ur{رہی} \ur{تھی۔}}
\end{quote}

\begin{quote}
\emph{bāhar bārish ho rahī thī.}
\end{quote}

\begin{quote}
\textbf{\ur{جب} \ur{گاڑی} \ur{لاہور} \ur{پہنچی} \ur{تو} \ur{دوست} \ur{وہاں} \ur{تھی۔}}
\end{quote}

\begin{quote}
\emph{jab gāṛī Lāhaur pahuṅchī to dost vahāṅ thī.}
\end{quote}

\begin{quote}
\textbf{\ur{دونوں} \ur{نے} \ur{اسٹیشن} \ur{پر} \ur{چائے} \ur{پی۔}}
\end{quote}

\begin{quote}
\emph{donoṅ ne isṭeshan par chāy pī.}
\end{quote}

\begin{quote}
\textbf{\ur{چائے} \ur{پی} \ur{کر} \ur{وہ} \ur{گھر} \ur{گئیں۔}}
\end{quote}

\begin{quote}
\emph{chāy pī kar voh ghar gaīṅ.}
\end{quote}

\begin{quote}
\textbf{\ur{شام} \ur{کو} \ur{انہوں} \ur{نے} \ur{کھانا} \ur{کھایا۔}}
\end{quote}

\begin{quote}
\emph{shām ko unhoṅ ne khānā khāyā.}
\end{quote}

\begin{quote}
\textbf{\ur{کھانے} \ur{کے} \ur{بعد} \ur{وہ} \ur{بازار} \ur{گئیں۔}}
\end{quote}

\begin{quote}
\emph{khāne ke bād voh bāzār gaīṅ.}
\end{quote}

\begin{quote}
\textbf{\ur{وہاں} \ur{سارا} \ur{نے} \ur{امی} \ur{کے} \ur{لیے} \ur{تحفہ} \ur{خریدا۔}}
\end{quote}

\begin{quote}
\emph{vahāṅ Sārā ne ammī ke liye tohfā kharīdā.}
\end{quote}

\begin{quote}
\textbf{\ur{آخر} \ur{میں} \ur{وہ} \ur{تھک} \ur{کر} \ur{سو} \ur{گئی۔}}
\end{quote}

\begin{quote}
\emph{ākhir meṅ voh thak kar so gaī.}
\end{quote}

\begin{quote}
\textbf{\ur{وہ} \ur{ایک} \ur{اچھا} \ur{دن} \ur{تھا۔}}
\end{quote}

\begin{quote}
\emph{voh ek achchhā din thā.}
\end{quote}

Now with nothing under it. Read it once without stopping:

\begin{quote}
\ur{پچھلے} \ur{ہفتے} \ur{سارا} \ur{لاہور} \ur{گئی۔} \ur{اس} \ur{کی} \ur{دوست} \ur{لاہور} \ur{میں} \ur{رہتی} \ur{ہے۔} \ur{پہلے} \ur{اس} \ur{نے} \ur{اسٹیشن} \ur{پر} \ur{ٹکٹ} \ur{لیا۔} \ur{پھر} \ur{وہ} \ur{ریل} \ur{گاڑی} \ur{میں} \ur{بیٹھی۔} \ur{سفر} \ur{کے} \ur{دوران} \ur{وہ} \ur{کتاب} \ur{پڑھ} \ur{رہی} \ur{تھی۔} \ur{باہر} \ur{بارش} \ur{ہو} \ur{رہی} \ur{تھی۔} \ur{جب} \ur{گاڑی} \ur{لاہور} \ur{پہنچی} \ur{تو} \ur{دوست} \ur{وہاں} \ur{تھی۔} \ur{دونوں} \ur{نے} \ur{اسٹیشن} \ur{پر} \ur{چائے} \ur{پی۔} \ur{چائے} \ur{پی} \ur{کر} \ur{وہ} \ur{گھر} \ur{گئیں۔} \ur{شام} \ur{کو} \ur{انہوں} \ur{نے} \ur{کھانا} \ur{کھایا۔} \ur{کھانے} \ur{کے} \ur{بعد} \ur{وہ} \ur{بازار} \ur{گئیں۔} \ur{وہاں} \ur{سارا} \ur{نے} \ur{امی} \ur{کے} \ur{لیے} \ur{تحفہ} \ur{خریدا۔} \ur{آخر} \ur{میں} \ur{وہ} \ur{تھک} \ur{کر} \ur{سو} \ur{گئی۔} \ur{وہ} \ur{ایک} \ur{اچھا} \ur{دن} \ur{تھا۔}
\end{quote}
\end{tcolorbox}

\begin{culture}
\textbf{\ur{گئی}} and \textbf{\ur{بیٹھی}} follow Sara: no object. \textbf{\ur{ٹکٹ} \ur{لیا}} and \textbf{\ur{تحفہ} \ur{خریدا}} follow the thing, both masculine; \textbf{\ur{چائے} \ur{پی}} follows \textbf{\ur{چائے}}, feminine. \textbf{\ur{پڑھ} \ur{رہی} \ur{تھی}} and \textbf{\ur{ہو} \ur{رہی} \ur{تھی}} are the scene.
\end{culture}

\subsection*{Your turn}
Say these aloud:

\begin{itemize}
\raggedright
  \item what Sara did first, at the station
  \item how the day ended, with \emph{ākhir meṅ}
\end{itemize}

\begin{tcolorbox}[breakable,colback=teal!4,colframe=teal!35!black,title={Before you move on}]
Why is \textbf{\ur{پڑھ} \ur{رہی} \ur{تھی}} in the scene form? (The reading was going on during the journey.) Which words keep the steps in order? (\textbf{\ur{پہلے}}, \textbf{\ur{پھر}}, \textbf{\ur{کے} \ur{بعد}}, \textbf{\ur{آخر} \ur{میں}}.)
\end{tcolorbox}
\section[(review)]{Review — a story in order}

\label{lesson:UR-R267-a-story-in-order}



\begin{quote}
Say \textquotedblleft{}I was reading\textquotedblright{} and \textquotedblleft{}during the meal\textquotedblright{}. (\emph{maiṅ paṛh rahā thā}, \emph{khāne ke daurān}.)
\end{quote}

\subsection*{Your turn}
\begin{itemize}
\raggedright
  \item \emph{Say it:} First I had a bath. Then I ate.
  \item \emph{Say it:} After that I went to the office.
  \item \emph{Say it:} I finished the work and came home.
  \item \emph{Say it:} In the end I drank tea and slept.
  \item \emph{Say it:} Sara's trip, in four lines: Lahore, the train, her friend, the market.
  \item \emph{Write it:} \textbf{\ur{بعد}} once, and point at \ur{ع} inside it
\end{itemize}

\begin{tcolorbox}[breakable,colback=teal!4,colframe=teal!35!black,title={Before you move on}]
\textquotedblleft{}After the meal\textquotedblright{}? (\textbf{\ur{کھانے} \ur{کے} \ur{بعد}}.) \textquotedblleft{}Having eaten\textquotedblright{}? (\textbf{\ur{کھا} \ur{کر}}.) Now the other way round: in \textbf{\ur{اس} \ur{نے} \ur{کھا} \ur{کر} \ur{چائے} \ur{پی}}, which verb decides \textbf{\ur{نے}}? (\textbf{\ur{پی}}, the last.) And the four words that put a story in order? (\textbf{\ur{پہلے}}, \textbf{\ur{پھر}}, \textbf{\ur{اس} \ur{کے} \ur{بعد}}, \textbf{\ur{آخر} \ur{میں}}.)
\end{tcolorbox}
